\section{Cook--Levin for \texorpdfstring{\(\mathcal L\)}{L}-descriptions}
\label{sec:l-cook-levin}

\begin{reminder}
A machine run is a path of configurations under one local transition rule;
Section~\ref{sec:refresher-tm} introduced this viewpoint.  Here the same idea
is applied to the fixed universal evaluator, turning a fuel-bounded lambda run
into locally checkable tableau data.
\end{reminder}

\subsection{A local encoding of bounded evaluation}

For closed \(t\), input \(u\), and fuel \(F\), the \emph{bounded trace} is
\[
 \mathcal T(t,u,F)=(C_0,C_1,\ldots,C_r),\qquad r\leq F,
\]
where \(C_0\) is the initial CEK configuration and each \(C_{i+1}\) is its
unique small-step successor.  A row records the current term pointer,
environment-frame pointers, continuation, heap, primitive microstate, outcome
flag, and binary fuel.  If evaluation halts early, a fixed self-looping final
state pads the trace to \(F+1\) rows.  This is the promised sequence of
\((\text{term},\text{environment},\text{fuel})\) configurations, with the
continuation and heap made explicit so that one-step validity is local.

To obtain a finite-alphabet tableau, compile the CEK evaluator and each charged
primitive to a fixed multitape Turing machine \(\widehat U_{\mathcal L}\).
The read-only code track contains \(\operatorname{ser}(t)\); input tracks hold
lazy argument tapes; work tracks hold the environment, heap, stack, and fuel.
Pointer arithmetic, traversal, and a primitive of charge \(k\) are expanded
into their elementary head moves.  One microstep changes the finite control,
the cells under the heads, and the head positions by at most one.

\begin{lemma}[Transition-window lemma]
\label{lem:transition-window}
There is a finite alphabet \(\Gamma_{\mathcal L}\), a fixed number \(r_0\) of
tracks, and a Boolean function
\[
 \tau:\bigl(\Gamma_{\mathcal L}^{3}\bigr)^{r_0}
       \times Q_{\mathcal L}\longrightarrow
       \Gamma_{\mathcal L}^{r_0}\times Q_{\mathcal L}
\]
with the following property.  For every nonboundary cell \(j\), the symbols in
cell \(j\) of row \(i+1\), together with the next finite-control state, are
determined by the radius-one windows
\((j-1,j,j+1)\) of row \(i\) and the current control state.  Boundary cells use
a fixed endmarker version of the same rule.  Conversely, if two consecutive
rows have one head marker on each track, agree with \(\tau\) at every window,
and have valid endmarkers, then the second row is exactly the successor of the
first under the small-step semantics.
\end{lemma}

\begin{proof}
Encode a head by pairing the tape symbol with one of the finitely many machine
states when that head is the active head, and use a separate finite-control
track for the other heads.  In one transition, a cell farther than one position
from every head is copied.  At a head cell, the transition table determines the
written symbol and next state from the old state and scanned symbols.  At its
left and right neighbors, the table determines whether the head marker enters;
it cannot enter any other cell because heads move by at most one.  These cases
depend only on the three displayed cells on each track.  They are mutually
exclusive once the old row has one head per track, and therefore define the
finite function \(\tau\).  Endmarkers replace missing neighbors by a fixed
boundary symbol.

For the converse, consider each track.  Away from the unique old head, the
local copy case forces equality.  In the three-cell neighborhood of that head,
the transition case forces the unique symbol write, direction, and new head
position specified by the machine table.  The local rules also force all other
cells not to acquire a head.  Repeating this argument for every track and
using the shared finite-control component gives precisely one transition of
\(\widehat U_{\mathcal L}\).  Primitive execution and CEK operations introduce
no exceptional case because they were compiled into the same elementary
machine steps.  The correspondence between those microstates and the weighted
semantics was part of the construction in Section~\ref{sec:resource-lambda}.
\end{proof}

\subsection{The succinct 3SAT circuit}

Let the row width be \(W\).  In \(F\) microsteps no head visits beyond the
initial code and input plus \(F\) cells, so
\(W=O(|t|+|u|+F)\).  Introduce Boolean variables for the constant-size binary
encoding of every tableau cell and control-state bit.  The formula asserts:
the initial row contains \(t,u,F\); every row and track has the prescribed
format and head marker; each adjacent pair obeys every transition window; and
the last row is accepting.  Each assertion involves a constant window except
binary addressing and the initial read-only strings.  Fixed Tseitin gadgets
turn it into 3CNF with \(O(FW)\) variables and clauses.

For the paper, the next theorem buys a small circuit describing an enormous
bounded verifier tableau, which is the input expected by answer reduction.
\begin{theorem}[Succinct Cook--Levin for \(\mathcal L\)]
\label{thm:l-succinct-sat}
There is an algorithm which, on
\((t,n,F,Q,\sigma,x,y)\), where \(|t|\leq\sigma\) and \(|x|,|y|\leq Q\),
outputs a Boolean circuit \(C_{t,n,F,Q,\sigma,x,y}\) deciding membership of an
indexed signed clause in a 3CNF formula \(\Phi\).  Its index width is
\(O(\log F+\log\sigma)\), its size and construction time are
\[
 \poly(\log n,\log F,Q,\sigma),
\tag{11.1}\label{eq:l-succinct-size}
\]
and for every pair of answer strings \(a,b\) the formula is satisfiable with
the prescribed input bits if and only if
\(\Eval_{\mathcal L}(\Quote(t),(n,x,y,a,b);F)\) accepts.  Holding the public
input fixed, the transition component of the circuit has size
\(\poly(\log F,|t|)\); the \(Q\) term in \eqref{eq:l-succinct-size} is solely
the hardwired initialization data \(x,y\).
\end{theorem}

\begin{proof}
Given a clause index, the circuit first decodes whether the requested clause is
an initialization, uniqueness/format, transition, or acceptance clause.  Row
and column numbers use \(O(\log F+\log W)\) bits.  For a transition clause it
computes the addresses of a radius-one window and evaluates the finite truth
table \(\tau\) from Lemma~\ref{lem:transition-window}; binary addition,
comparison, and multiplexing have size polynomial in those address lengths.
For an initialization clause it selects a bit of the canonical code or public
input.  A balanced selector for a hardwired string has size linear in
\(|t|+Q+\log n\).  Format and final-state clauses use fixed truth tables.
Finally it emits the three variable addresses and signs of the selected
3CNF gadget.  This proves the circuit and construction bounds.  Since
\(W=O(|t|+Q+F)\), \(\log W=O(\log F+\log\sigma+\log(Q+1))\), which is absorbed
by \eqref{eq:l-succinct-size} under \(Q\leq F\), as in the paper.

If evaluation accepts, assign every tableau variable from its unique bounded
trace and every Tseitin variable from its gadget value.  Initialization,
transition, and acceptance clauses are then satisfied.  Conversely, any
satisfying assignment contains the prescribed initial row.  Induction on the
row number and the converse part of Lemma~\ref{lem:transition-window} forces
each later row to be the unique evaluator successor.  The final clauses force
that trace to accept.  The prescribed answer bits occupy the initial answer
tapes, so the equivalence holds for each \(a,b\).
\end{proof}

This is the \(\mathcal L\)-description restatement of
\code{prop:standard-succinct-sat}.  The paper's version takes \(D,n,T,Q,
\sigma,x,y\), reserves the first \(2T\) encoded positions for each answer, and
has size \(\poly(\log n,\log T,Q,\sigma)\).
\gt{prop:standard-succinct-sat}{ground-truth/gt-10-answer-reduction.tex}
Theorem~\ref{thm:l-to-tm} gives the same proposition by compilation; the proof
above shows directly why a lambda evaluation trace has the required locality.

More explicitly, use the paper's tape-alphabet encoding
\(\operatorname{enc}_\Gamma(0)=00\),
\(\operatorname{enc}_\Gamma(1)=01\),
\(\operatorname{enc}_\Gamma(\sqcup)=10\), and reserve \(2F\) bits in the
tableau assignment for each answer tape.  For every
\(a,b\in\{0,1\}^{2F}\), the resulting formula has an extension \(c\) if and
only if there are prefixes \(a',b'\), of lengths at most \(F\), such that
\[
 a=\operatorname{enc}_\Gamma(a',\sqcup^{F-|a'|}),\qquad
 b=\operatorname{enc}_\Gamma(b',\sqcup^{F-|b'|}),
\]
and the term accepts \((n,x,y,a',b')\) within fuel \(F\).  The forward
direction follows because the initialization clauses force the two reserved
blocks to be valid padded tape rows; the backward direction fills those blocks
from the accepting trace.  This is the exact placement property needed when
the two blocks are separated in \eqref{eq:decoupled-five}.

\subsection{From shared 3SAT data to decoupled 5SAT}

The paper removes two undesirable features before low-degree encoding.  A
3SAT clause reads three locations of one tableau assignment, while three
low-degree answers need not be restrictions of one assignment.  Also, the
strings \(a,b\) are embedded inside that assignment, whereas the later sampler
must query them as independent blocks.  Make five assignments
\(a,b,u_3,u_4,u_5\) and impose clauses of the form
\[
 a_{i_1}^{o_1}\vee b_{i_2}^{o_2}\vee
 (u_3)_{i_3}^{o_3}\vee (u_4)_{i_4}^{o_4}\vee
 (u_5)_{i_5}^{o_5}.
\tag{11.2}\label{eq:decoupled-five}
\]
The original 3SAT clauses are applied to \(u_3\); constant-size five-literal
gadgets enforce \(u_3=u_4=u_5\), copy the first \(2F\) encoded answer bits to
\(a,b\), and pad unused positions by the alternating blank encoding.  Each
gadget is indexed by comparisons and additions on
\(O(\log F+\log\sigma)\)-bit addresses, so it adds only that many gates.
Therefore the succinct decoupled circuit retains the bound
\eqref{eq:l-succinct-size}.  Its padded form makes all five blocks have one
common power-of-two length.  A satisfying five-tuple restricts to a satisfying
trace assignment, and a satisfying trace extends by copying and padding; this
proves the required if-and-only-if.
\gt{sec:succinct-deciders}{ground-truth/gt-10-answer-reduction.tex}

The intended analytic chain is
\[
\begin{split}
\mathtt{Quoted\{Decider\}}
&\xrightarrow{\ \mathtt{bounded\_trace}\ }
 \mathtt{BoundedTrace}\\
\mathtt{BoundedTrace}
&\xrightarrow{\ \mathtt{cook\_levin}\ }
 \mathtt{Succinct3SAT}\\
\mathtt{Succinct3SAT}
&\xrightarrow{\ \mathtt{decouple5}\ }
 \mathtt{SuccinctDecoupled5SAT}
\\
\mathtt{SuccinctDecoupled5SAT}
&\xrightarrow{\ \mathrm{padding}\ }
 \mathtt{PaddedSuccinctDecoupled5SAT}.
\end{split}
\]
These are the transformations planned for TB3 in \code{briefs/23-tb3.md}.
That brief is planned work, not evidence that the functions exist or that any
claim has been checked.  At the time of this note, no \code{bounded_trace},
\code{cook_levin}, or \code{decouple5} implementation is present in
\code{src/}; the general contracts remain \grade{CITED}.

\subsection{Tseitin and arithmetization}

\begin{reminder}
At this point the quoted verifier has already been reduced to a finite circuit
describing local computation checks.  The transformations below manipulate
that circuit description; they are computable source maps, while their
low-degree and quantum soundness consequences are separate analytic claims.
\end{reminder}

Given a Boolean circuit \(C\) with gates \(g_i\), Tseitin introduces a wire
variable \(w_i\) and an equivalence formula
\[
 z_i=(g_i(x,w)\wedge w_i)\vee
     (\neg g_i(x,w)\wedge\neg w_i),
 \qquad F=\bigwedge_i z_i.
\]
The required contract is
\(C(x)=1\) iff there exists \(w\) with \(F(x,w)=1\), with size and
construction time \(O(|C|)\).\gt{def:tseitin}{ground-truth/gt-10-answer-reduction.tex}
The project follows the explicit gate-equivalence form in
\code{ground-truth/nw19/nw19-tseitin-arith.tex} and also conjoins the output
literal \(w_{\rm out}\), so the stored formula explicitly asserts acceptance.
That convention is marked \grade{SOURCE\_REPAIR}; it is not silently attributed
to the source.

The Julia constructors \code{Circuit}, \code{Input}, \code{Gate},
\code{NotGate}, \code{AndGate}, and \code{OrGate} make the finite DAG and its
topological order explicit.  The function \code{tseitin} returns
\code{Checked{TseitinFormula,CertNode}}.  Its replay recomputes the formula's
occurrence vector; it does not, by that replay alone, establish the general
Cook--Levin or Tseitin semantic equivalence.  Those implications are exhausted
only on the small TB0 fixture and remain cited in general.

Arithmetization recursively applies
\[
  \neg A\mapsto 1-A,\qquad
  A\wedge B\mapsto AB,\qquad
  A\vee B\mapsto A+B-AB.
\]
It agrees with the Boolean formula on the Boolean cube.
\gt{def:formula-arithmetization}{ground-truth/gt-10-answer-reduction.tex}
The related NW19 rule is \code{def:arithmetization} in
\code{ground-truth/nw19/nw19-tseitin-arith.tex}.  Julia's \code{arith_q}
returns a sparse formal \code{Poly}; its \code{ArithTseitin} certificate checks
support degrees against a structural bound.

\subsection{The occurrence-vector discrepancy}

There is a discrepancy in the individual-degree bookkeeping as we read NW19;
our reading may be wrong.  The source proposition states individual degree at
most two for the Tseitin arithmetization.
\gt{prop:tseitin-arith-degree}{ground-truth/gt-10-answer-reduction.tex}
Reading the NW19 formula literally as a formula tree, arithmetization is
multilinear in leaf occurrences, not necessarily in circuit wires.  The same
wire can occur in several gate gadgets.  With the project's explicit output
literal, the count we compute is
\[
 \operatorname{occ}(x_j)=2\operatorname{fanout}(x_j),\qquad
 \operatorname{occ}(w_i)=2+2\operatorname{fanout}(w_i)
       +\mathbf 1_{i=\mathrm{out}},
\]
and structural induction gives
\[
 \deg_v(F_{\rm arith})\leq \operatorname{occ}_v(F).
\]
The function \code{tseitin_occurrence_account} computes this vector;
\code{occurrences} independently traverses the formula; and \code{arith_q}
compares the structural account with actual sparse support.

Claim C8 is only \claim{C8}{TESTED}: the two-gate regression gives
\((2,2,2,4,3)\), and the 16-variable TB0 formula attains
\[
(2,0,0,0,2,2,0,0,0,0,6,4,4,4,4,3).
\]
The general occurrence formula is a written derivation, not a machine theorem.
The downstream soundness repair is only \claim{C5}{SKETCH}.  With a symbolic
bound \(\delta_F=\max_v\operatorname{occ}_v(F)\), the conservative formula-test
term becomes \((\delta_F+5d)m'/q\), and
\code{P_formula_structural} is an additional obligation.  Under an explicit
fanout-reduction hypothesis, \(\operatorname{fanout}_{\max}\leq2\) gives
\(\delta_F\leq7\) with the output literal, and \(d\geq8\) suffices for the
individual-degree construction bound.  Absorption into the asymptotic
parameter choice still needs the stated growth hypotheses; it is not promoted
by the two finite computations.

\section{The low-degree PCP as algebra}

\subsection{Interpolation and the vanishing polynomial}

For \(M=2^m\) and \(a=(a_y)_{y\in\{0,1\}^m}\), define
\[
 \operatorname{ind}_{m,y}(x)=
 \prod_{j:y_j=1}x_j\prod_{j:y_j=0}(1-x_j),\qquad
 g_a(x)=\sum_y a_y\operatorname{ind}_{m,y}(x).
\]
Then \(g_a(y)=a_y\) on the Boolean cube and \(a\mapsto g_a\) is linear.
This is the low-degree code: the assignment is represented by the evaluation
table of its unique multilinear extension.
\gt{sec:ld-encoding}{ground-truth/gt-03-prelim.tex}
The Julia function \code{g_a} constructs the sparse polynomial; block
coordinates are carried by \code{VarLayout} and \code{VarBlock};
\code{dependency_blocks} recomputes the actual support dependency.

For five satisfying blocks, let \(g_i\) be their multilinear extensions and
write \(z=(x_1,\ldots,x_5,o,w)\in\F_q^{m'}\).  The central analytic object is
\[
 c_0(z)=F_{\rm arith}(z)
        \prod_{i=1}^{5}\bigl(g_i(x_i)-o_i\bigr).
\]
On a Boolean point, either \(F_{\rm arith}=0\), or the indicated clause is
present and at least one factor \(g_i(x_i)-o_i\) vanishes.  Hence \(c_0\)
vanishes on \(\{0,1\}^{m'}\).  The code executes the product in
\code{build_c0}; the \code{BuildC0} certificate replays the structural and
actual degree accounts.  The implication from a satisfying decoupled instance
to Boolean-cube vanishing is checked only on the TB0 fixture and cited in
general.\gt{thm:pcp-decider}{ground-truth/gt-10-answer-reduction.tex}

\subsection{Division by the Boolean ideal}

Let \(z_i(1-z_i)=z_i-z_i^2\).  The ground-truth basis proposition says that a
degree-bounded polynomial vanishing on the Boolean cube lies in the ideal
generated by these \(m'\) polynomials.
\gt{prop:zero-basis}{ground-truth/gt-10-answer-reduction.tex}
The implementation uses the following deterministic monomial rewrite.

\begin{lemma}[Executable zero-basis rewrite]
For a coefficient \(a\), a monomial \(u\) independent of the displayed power,
and \(e\geq2\),
\[
 a u z_i^e
 \longmapsto
 a u z_i^{e-1}-a u z_i^{e-2}\,z_i(1-z_i).
\]
Iterating in a fixed variable and monomial order yields polynomials \(c_i\) and
a remainder \(r\), multilinear in every variable, such that
\[
 c_0=\sum_{i=1}^{m'}c_i z_i(1-z_i)+r.
\]
If the coordinate degree vector of \(c_0\) is \(D\), then
\(\deg_j(c_i)\leq D_j\), \(\deg_i(c_i)\leq\max(D_i-2,0)\), and after reducing
coordinates \(1,\ldots,i-1\), the running remainder has degree at most one in
those coordinates.
\end{lemma}

\begin{proof}
Expanding the right-hand side of one rewrite gives
\[
 a u z_i^{e-1}-a u z_i^{e-2}(z_i-z_i^2)=a u z_i^e.
\]
Thus every step preserves the coefficient polynomial after adding the recorded
quotient contribution.  It lowers the current exponent to one and places
power \(e-2\) in the quotient, proving the degree bounds.  Induction over the
ordered coordinates gives the displayed decomposition and a final
multilinear remainder.  If \(c_0\) vanishes on the Boolean cube, then so does
\(r\); uniqueness of multilinear interpolation on \(\{0,1\}^{m'}\) makes
\(r=0\).  This is the constructive division used in the proof of
\code{prop:zero-basis} in
\code{ground-truth/gt-10-answer-reduction.tex}.
\end{proof}

\code{zero_basis_decompose} stores every step as \code{RewriteStep} inside a
\code{ZeroDecomposition}.  \code{verify_rewrite_step} replays the scalar
identity, and \code{verify_zero_decomposition} reconstructs the entire right
side as a normalized coefficient dictionary.  Promotion to a zero certificate
also requires empty remainder.  This is stronger than checking the equality at
sampled points, but it remains an instance certificate, not a proof of the PCP
soundness theorem.

\subsection{PCP view and local verifier}

A \emph{PCP view} at \(z\) is
\[
 \operatorname{ev}_z(\Pi)=
 (g_1(x_1),\ldots,g_5(x_5),c_0(z),c_1(z),\ldots,c_{m'}(z)).
\]
The proof object is the tuple of evaluation tables for those polynomials.
The labels are \code{def:pcp-proof} and \code{def:pcp-eval} in
\code{ground-truth/gt-10-answer-reduction.tex}.  Julia represents them by
\code{PCPProof} and \code{PCPView}; \code{ev_z} refuses a \(g_i\) whose actual
support escapes block \(X_i\).

The verifier first reconstructs the succinct circuit, Tseitin formula, and
arithmetization.  It then checks exactly
\[
 \beta_0=F_{\rm arith}(z)\prod_{i=1}^5(\alpha_i-o_i),
 \qquad
 \beta_0=\sum_{j=1}^{m'}\beta_j z_j(1-z_j).
\]
These are the formula test and zero-on-subcube test.
\gt{fig:pcpverifier}{ground-truth/gt-10-answer-reduction.tex}
The Julia function \code{pcpverifier} executes these two equations on an
already supplied \code{TseitinFormula} and view.  In TB2 it does not implement
the figure's preceding reconstruction steps; the formula is construction-time
data.  \code{build_pcp} combines upstream certificates, coefficient division,
degree checking, and stored certified views.

\subsection{The four soundness layers}

The four layers must not be collapsed.

\paragraph{1. Algebraic soundness under a low-degree premise.}
Assume every submitted \(g_i,c_j\) has individual degree at most \(d\).  If the
formula-test polynomials are unequal, their difference has total degree at most
\((\delta_F+5d)m'\) under the occurrence-vector account.  If the zero-test
polynomials are unequal, their difference has total degree at most
\((2+d)m'\).  Schwartz--Zippel states that unequal polynomials of total degree
at most \(\Delta\) agree at a uniform point of \(\F_q^{m'}\) with probability
at most \(\Delta/q\).
\gt{lem:schwartz-zippel}{ground-truth/gt-03-prelim.tex}
Consequently
\[
 \Pr[\text{false formula identity}]
 \leq \frac{(\delta_F+5d)m'}{q},\qquad
 \Pr[\text{false zero identity}]
 \leq \frac{(2+d)m'}{q}.
\]
The paper's displayed first numerator uses \(2+5d\); the structural project
obligation substitutes the measured \(\delta_F\) under the discrepancy stated
above.  The field size enters nowhere mysteriously: \(q=2^k\) is chosen so
both fractions are below the acceptance threshold \(1/2\), together with the
other PCP parameter obligations.  The first structural inequality is extra to
the literal \code{def:pcpparams} predicate in the present project.

If acceptance is greater than \(1/2\) and both upper bounds are below
\(1/2\), each sampled equality is therefore a polynomial identity.  The zero
identity makes \(c_0\) vanish on the Boolean cube; the formula identity then
decodes five Boolean assignments satisfying the decoupled formula, and hence an
accepting bounded computation.  This conditional theorem is
\code{thm:pcp-decider} in
\code{ground-truth/gt-10-answer-reduction.tex}.  Its current project status is
\claim{C5}{SKETCH}, not CHECKED.

\paragraph{2. Enforcement of the low-degree premise.}
The simultaneous low-degree game queries points, axis-parallel lines, and
diagonal lines.  Its classical decider checks consistency of a point value with
a claimed line restriction.\gt{fig:ld-decider}{ground-truth/gt-07-ldt.tex}
The theorem extracting approximately consistent low-individual-degree
polynomial measurements from a successful quantum strategy is
\code{lem:ld-soundness} in \code{ground-truth/gt-07-ldt.tex}.  It is
\grade{CITED}.  Running \code{ld_decider} on honest finite inputs does not prove
this extraction theorem.

\paragraph{3. Polynomial identity from random evaluation.}
This is exactly the Schwartz--Zippel step above, used twice with two different
degree bounds.  The local program checks evaluations.  The analytic theorem
turns an acceptance probability into formal identity, provided the degree and
field-size premises hold.  No sample count can replace that implication.

\paragraph{4. Quantum consistency and rigidity.}
Answer reduction must connect different players' point and line answers to the
same polynomial measurements, control disturbance, and transport the result
back to the original verifier.  Those operator estimates use low-degree
soundness, typed consistency transformations, and later rigidity arguments.  They
are imported in the soundness contract of \code{thm:ar} in
\code{ground-truth/gt-10-answer-reduction.tex}; none is a sparse-polynomial
identity.

The computer adds exactly this: on a named instance it constructs the
polynomials, executes the division, compares normalized coefficients, computes
degree and dependency vectors, and evaluates the two local equations.  It does
not test the theorem.  Current statuses reflect that separation: C3 is
\textsc{TESTED}, C5 is \textsc{SKETCH}, C1 and C2 remain
\textsc{CONJECTURE}, and C8 is \textsc{TESTED} only on its two named circuits.

\section{The typed layer: conditionally linear functions}

\subsection{The inductive object and its laws}

\begin{definition}[Conditionally linear function]
Let \(V=\F^n\) with fixed coordinate-register subspaces.  The unique level-zero
conditionally linear (CL) function is zero.  A level-\(\ell\) CL function
chooses complementary registers \(V_1\oplus V_{>1}=V\), applies a linear map
\(A:V_1\to V_1\), and, conditional only on \(A(x^{V_1})\), selects a
level-\((\ell-1)\) CL function on \(V_{>1}\).  Its output is the sum of the
first-stage and continuation outputs.
\end{definition}
This is \code{def:cl-func} in
\code{ground-truth/gt-04-cl.tex}.  Julia mirrors it with
\code{AbstractCL}, \code{CLZero}, and \code{CLStep}.  The constructor verifies
disjoint factor/rest coordinate sets, a matrix acting inside the factor, a
fixed child level, and exact coverage of the rest register.  Thus an inhabitant
of \code{CLStep} has a constructed level; linearity is carried by the matrix,
not inferred from samples.

The marginal \(L_{\leq k}\) is the sum of the first \(k\) stage outputs,
together with their factor spaces and linear maps.  This is the decomposition
of \code{lem:cl-kth} in
\code{ground-truth/gt-04-cl.tex}.  Julia's \code{marginal_k} returns a
\code{CLMarginal}; \code{sum_stage_outputs} recomputes its value.

\begin{lemma}[Level laws]
Let \(L\) be level \(k\), let every continuation \(R_u\) be level \(\ell\),
and let \(L_j\) act on pairwise disjoint registers at levels \(\ell_j\).
Then
\[
 \operatorname{level}(\operatorname{concatenate}(L,R))=k+\ell,
 \qquad
 \operatorname{level}\!\left(\bigoplus_jL_j\right)=\max_j\ell_j.
\]
For CL distributions \(S_1,S_2\), their product, formed by direct-summing the
left maps and the right maps on independent seed registers, has level the
maximum of the two levels.
\end{lemma}

\begin{proof}
For concatenation, induct on the outer constructor of \(L\).  At level zero,
grafting \(R_0\) supplies exactly \(\ell\) stages.  At a \code{CLStep}, retain
its first linear stage and apply the induction hypothesis to each child; this
gives \(1+(k-1+\ell)=k+\ell\).  This is the construction in
\code{lem:cl-concat} of \code{ground-truth/gt-04-cl.tex}.

For direct sum, pad each component by zero stages to
\(r=\max_j\ell_j\).  At stage \(i\), use the direct sum of the components'
stage matrices on the direct sum of their factor registers.  The continuation
depends only on the tuple of earlier outputs, hence remains CL.  Induction on
\(r\) proves the law, matching \code{lem:cl-func-prod} in
\code{ground-truth/gt-04-cl.tex}.  A distribution product applies this law
separately to its left and right maps, so its level is the same maximum.
\end{proof}

The code functions are \code{concatenate}, \code{direct_sum},
\code{distribution}, and \code{product}.  Their type constructions carry the
level laws.  The stronger structural hypothesis that all compiler operations
of interest close inside one useful algebra is not implied by these individual
lemmas.

\subsection{Point and line maps}

For \(V=V_{\rm pt}\oplus V_{\rm coord}\oplus V_{\rm dir}\), write a seed as
\((u,s,v)\in\F_q^m\times\F_q\times\F_q^m\).  Let
\(L_v^{\rm lnf}\) be the canonical projection whose kernel is
\(\operatorname{span}(v)\), so \(L_v^{\rm lnf}u\) is a canonical representative
of the affine line \(u+\F_qv\).
\gt{def:cl-canonical}{ground-truth/gt-03-prelim.tex}
The line-representative use is \code{def:line-representative} in
\code{ground-truth/gt-07-ldt.tex}.  Define \(\chi(s)\) by splitting the fixed
integer representatives of \(\F_q\) into \(m\) equal buckets; this requires
\(m\mid q\).  With \(\pi_{i-1}\) zeroing the first \(i-1\) coordinates,
\[
\begin{aligned}
L_{\rm Point}(u,s,v)&=(u,0,0),\\
L_{\rm ALine}(u,s,v)&=(L_{e_{\chi(s)}}^{\rm lnf}u,s,0),\\
L_{\rm DLine}(u,s,v)&=(L_{\pi_{\chi(s)-1}(v)}^{\rm lnf}u,
                         s,\pi_{\chi(s)-1}(v)).
\end{aligned}
\]
Their levels are respectively one, two, and three by their stage
factorizations.\gt{sec:ld-game}{ground-truth/gt-07-ldt.tex}
The implementation names are \code{L_Point}, \code{L_ALine},
\code{L_DLine}, \code{L_lnf}, \code{chi}, and \code{pi_prefix}.  Since the
source's kernel-basis definition does not cover \(v=0\), the code uses
\(L_0^{\rm lnf}=I\) and marks the convention \grade{SOURCE\_REPAIR}.

The distribution \(\mu_{L,R}\) pushes one uniform shared seed through the two
maps.\gt{def:cl-dist}{ground-truth/gt-04-cl.tex}
For \((q,m)=(8,2)\), TB1 enumerates all \(8^5=32768\) seeds.  The
axis-line/point histogram has support 512, and the diagonal-line/point
histogram has support 18432, including the zero directions.  Both agree
entry-for-entry with separately transcribed formulas for
\code{lem:alnf} and \code{lem:dlnf} in
\code{ground-truth/gt-07-ldt.tex}; every marginal is replayed.  This finite
statement is \claim{C4a}{TESTED}.  It is not a statistical argument and does
not establish quantum low-degree soundness.

\subsection{Types and the six-copy sampler}

A \emph{typed sampler} is a finite type-graph-indexed family of CL maps sharing
one ambient seed space.  It first samples an oriented type edge and then pushes
one uniform seed through the selected left/right maps.
\gt{def:typed-sampler}{ground-truth/gt-06-types.tex}
Its induced distribution is \code{def:typed-sampler-sample} in the same file.
Julia's \code{TypedSampler} checks exact type coverage, a common field and seed
dimension, graph endpoints, and a padded common level.

\emph{Detyping} encodes the external graph choice into an untyped CL
distribution by a rejection-sampling construction.  It adds two levels,
preserves value-one completeness, and changes soundness error by
\(16^{|\mathrm{Type}|}\).\gt{lem:detyping-verifiers}{ground-truth/gt-06-types.tex}
The Julia function \code{detype} returns \code{CitedDetypedVerifier}; the
certificate is deliberately \grade{CITED}.  For the 54 answer-reduction types,
the displayed factor is \(16^{54}\), not an executed estimate.

The answer-reduction PCP sampler has the 18 types
\[
 \{\mathrm{Point}_i,\mathrm{ALine}_i,\mathrm{DLine}_i:1\leq i\leq6\}
\]
and a complete type graph.  Copies 1--5 query one \(m\)-variate \(g_i\) each;
copy 6 queries the bundle \((g_1,\ldots,g_5,c_0,\ldots,c_{m'})\) in dimension
\(m'\).  The six-copy construction and its ambient direct sums are in
\code{sec:ld-compiler} of
\code{ground-truth/gt-10-answer-reduction.tex}.  The current code uses
\code{PCPType}, \code{PCPRegisterLayout}, \code{PCPCLMap},
\code{PCPPaddedMap}, and \code{pcp_sampler}.  It carries a visible
\grade{SOURCE\_REPAIR} for the copy-6 coordinate-register conflict between the
displayed ambient space and the parser table.

\emph{Oracularization} is the typed transformation in which the isolated
roles receive the original left or right CL image while an oracle role receives
the common seed and can reconstruct both; the code name is
\code{oracularize_sampler}.  Its theorem contract is \grade{CITED}.
\gt{sec:orac-def}{ground-truth/gt-09-oracularization.tex}
The typed
answer-reduced sampler is the product of this three-role family with the
18-type PCP family.  \code{typed_sampler_product} produces 54 types and
combines levels by a maximum.  The code constructs the product and checks its
shape, but C4b remains \textsc{CONJECTURE}: the latest TB2 critic held promotion
because the lazy \code{PCPCLMap}'s stage depth was not itself enough to make
the full CL decomposition a constructor-level fact.  Local certificate labels
must not be confused with a promoted claim.

The typed decider implements the five guarded checks of
\code{fig:decider-pcp}: global consistency, input consistency, input low degree,
proof encoding (individual and simultaneous low degree), and the PCP game
check.\gt{fig:decider-pcp}{ground-truth/gt-10-answer-reduction.tex}
The Julia names are \code{TypedAnswerReducedDecider},
\code{typed_answer_reduced_decider}, \code{LDParams}, \code{ld_decider}, and
\code{PCPGameCall}.  \code{answer_reduce_pcp} constructs the typed verifier of
level \(\max(\ell,3)\) on its fixture and attaches explicit assumptions.

The full theorem contract is conditional:
\[
\begin{array}{ll}
\textsc{Assume:}&V\text{ is an }\ell\text{-level normal-form verifier},\\
&T(n)=(2^{\lambda n})^\mu,\quad Q(n)=(\lambda n)^\mu,\\
&\Time_D(n)\leq T(n),\quad \Time_S(n)\leq Q(n);\\[2pt]
\textsc{Prove:}&\AnswerReduce(V)\text{ has level }\max\{\ell+2,5\},\\
&\text{the stated polynomial runtimes and directed completeness,}\\
&\text{soundness, and entanglement implications.}
\end{array}
\]
This is \code{thm:ar} in
\code{ground-truth/gt-10-answer-reduction.tex}.  The finite guarded decider is
executable; the quantum implication, detyping, and asymptotic theorem are
\grade{CITED}.

\subsection{The structural hypothesis}

\emph{Introspection} is the cited verifier transformation that makes encoded
prover-side measurement information available to a smaller verifier while
preserving a directed soundness contract.
\gt{thm:introspection}{ground-truth/gt-08-introspection.tex}
\emph{Anchoring} inserts special anchor questions with positive probability;
\emph{parallel repetition} runs several copies with an AND acceptance
condition against possibly correlated strategies.  Their combined repetition
contract is cited from \code{thm:repetition} in
\code{ground-truth/gt-11-parallel-repetition.tex}.

The relevant closure equations are
\[
 \Introspect(\mathcal Q)\subseteq\mathcal Q,
 \qquad \mathsf{PCPCompose}(\mathcal Q)\subseteq\mathcal Q,
 \qquad \mathcal Q\times\mathcal Q\subseteq\mathcal Q.
\]
Concatenation, direct sum, the shared-seed distribution, and product are visibly
CL-preserving in the base inductive datatype.  Affine restriction is linear
after earlier stages choose a direction, and random-point identity tests are
pushforwards of a shared uniform seed.  Typed graph choice is external, and the
cited detyping transformation moves it back into two further CL levels.

This does not characterize every admissible query distribution, prove that a
generic PCPP can be expressed in the algebra, or prove closure under the full
introspection compiler.  The structural statement is exactly
\claim{C7}{CONJECTURE}.  Introspection soundness and normal-form closure are
not implemented.  The present datatype makes several local closure laws
legible, not the global compression theorem.

\section{Correspondence tables}

The evidence grade describes the local certificate leaf.  The claim column is
the current status in \code{claims/CLAIMS.md}; it can remain weaker than a
certificate label when an adversarial verdict has not promoted the claim.
\grade{CONSTRUCTED} means enforced by a constructor,
\grade{CHECKED} means deterministic replay against an attached object, and
\grade{CITED} means imported from the named ground-truth statement.

\subsection{Description front end (planned TB3)}

\setlength{\tabcolsep}{1.5pt}
\scriptsize
\begin{longtable}{P{0.21\textwidth}P{0.16\textwidth}P{0.19\textwidth}P{0.19\textwidth}P{0.12\textwidth}P{0.08\textwidth}}
\toprule
Paper object & Analytic statement & IR sort & Julia name & Grade & Claim \\
\midrule
\endhead
\code{sec:tms}\newline\code{gt-03-prelim.tex}
& Machine conventions in Definition~\ref{def:paper-tm}; costed bridge in
Theorems~\ref{thm:tm-to-l}--\ref{thm:l-to-tm}
& \code{Quoted{A}} & \code{Quote}, \code{Eval}, \code{Specialize}
& \grade{CITED}; absent & none \\
\code{def:decider}\newline\code{gt-05-games-normalform.tex}
& Five-input format in Definition~\ref{def:paper-decider}; bounds preserved by
Theorem~\ref{thm:lambda-preservation} & \code{Quoted{Decider}}
& \code{description_size} & \grade{CITED}; absent & none \\
\code{DESIGN} \S1.1
& Determinism, fuel, quotation, and specialization:
Theorems~\ref{thm:l-determinism}--\ref{thm:quote-eval} and
Lemma~\ref{lem:specialization}
& \code{ClosedProgram}, \code{Quoted} & \code{Eval}, \code{specialize}
& analytic only; absent & none \\
\code{prop:standard-succinct-sat}\newline\code{gt-10-answer-reduction.tex}
& Transition locality (Lemma~\ref{lem:transition-window}) and succinct 3SAT
(Theorem~\ref{thm:l-succinct-sat}) & \code{BoundedTrace}, \code{Succinct3SAT}
& \code{bounded_trace}, \code{cook_levin} & \grade{CITED}; absent & none \\
\code{sec:succinct-deciders}\newline\code{gt-10-answer-reduction.tex}
& Three shared reads become five independent blocks &
\code{SuccinctDecoupled5SAT} & \code{decouple5} & \grade{CITED}; absent & none \\
\bottomrule
\end{longtable}
\normalsize

\subsection{Polynomial rung (TB0)}

\scriptsize
\begin{longtable}{P{0.21\textwidth}P{0.16\textwidth}P{0.19\textwidth}P{0.19\textwidth}P{0.12\textwidth}P{0.08\textwidth}}
\toprule
Paper object & Analytic statement & IR sort & Julia name & Grade & Claim \\
\midrule
\endhead
\code{def:tseitin}\newline\code{gt-10-answer-reduction.tex}
& Circuit becomes a linear-size formula with auxiliary wires &
\code{Circuit}, \code{TseitinFormula} & \code{tseitin}, \code{occurrences}
& \grade{CHECKED} occurrence & C2 \textsc{Conj.} \\
\code{def:formula-arithmetization}\newline\code{gt-10-answer-reduction.tex}
& Boolean connectives become a formal polynomial & \code{Poly}
& \code{arith_q}, \code{actual_degrees} & \grade{CHECKED} & C8 \textsc{Tested} \\
\code{sec:ld-encoding}\newline\code{gt-03-prelim.tex}
& Assignment becomes its multilinear extension & \code{Poly}
& \code{g_a}, \code{dependency_blocks} & \constructed/\newline\grade{Checked}
& C3 \textsc{Tested} \\
\code{prop:zero-basis}\newline\code{gt-10-answer-reduction.tex}
& Divide by \(z_i(1-z_i)\) and require zero remainder &
\code{ZeroDecomposition} & \code{zero_basis_decompose},
\code{verify_zero_decomposition} & \grade{CHECKED} & C2 \textsc{Conj.} \\
\code{fig:pcpverifier}\newline\code{gt-10-answer-reduction.tex}
& Check formula and zero identities at one view & \code{PCPView}
& \code{pcpverifier} & \grade{CHECKED} finite views & C1 \textsc{Conj.} \\
\code{thm:pcp-decider}\newline\code{gt-10-answer-reduction.tex}
& Acceptance \(>1/2\) decodes an accepting trace, assuming low degree &
\code{PCPProof} & \code{build_pcp}, \code{parameter_policy}
& \grade{CITED}/\newline\grade{Checked} parts & C5 \textsc{Sketch} \\
\bottomrule
\end{longtable}
\normalsize

\subsection{CL and low-degree-test rung (TB1)}

\scriptsize
\begin{longtable}{P{0.21\textwidth}P{0.16\textwidth}P{0.19\textwidth}P{0.19\textwidth}P{0.12\textwidth}P{0.08\textwidth}}
\toprule
Paper object & Analytic statement & IR sort & Julia name & Grade & Claim \\
\midrule
\endhead
\code{def:cl-func}\newline\code{gt-04-cl.tex}
& CL functions are inductive conditional linear stages & \code{AbstractCL}
& \code{CLZero}, \code{CLStep}, \code{level} & \constructed
& C4a \textsc{Tested} \\
\code{lem:cl-kth}\newline\code{gt-04-cl.tex}
& A marginal is the sum of the first \(k\) stages & \code{CLMarginal}
& \code{marginal_k}, \code{sum_stage_outputs} & \grade{CHECKED}
& C4a \textsc{Tested} \\
\code{lem:cl-concat}, \code{lem:cl-func-prod}\newline\code{gt-04-cl.tex}
& Concatenation adds levels; direct sum takes a maximum & \code{AbstractCL}
& \code{concatenate}, \code{direct_sum}, \code{product}
& \constructed & C7 \textsc{Conj.} \\
\code{lem:alnf}, \code{lem:dlnf}\newline\code{gt-07-ldt.tex}
& CL pushforwards equal the line/point distributions & \code{CLDistribution}
& \code{L_Point}, \code{L_ALine}, \code{L_DLine}, \code{histogram}
& \grade{CHECKED} TB1 & C4a \textsc{Tested} \\
\code{fig:ld-decider}\newline\code{gt-07-ldt.tex}
& Point/line answer consistency and degree format & \code{LDParams}
& \code{ld_decider}, \code{restrict} & \grade{CHECKED} fixture & C4a \textsc{Tested} \\
\code{lem:ld-soundness}\newline\code{gt-07-ldt.tex}
& Successful quantum strategy yields polynomial measurements & theorem leaf
& no proof checker & \grade{CITED} & C5 \textsc{Sketch} \\
\bottomrule
\end{longtable}
\normalsize

\subsection{Typed answer-reduction rung (TB2)}

\scriptsize
\begin{longtable}{P{0.21\textwidth}P{0.16\textwidth}P{0.19\textwidth}P{0.19\textwidth}P{0.12\textwidth}P{0.08\textwidth}}
\toprule
Paper object & Analytic statement & IR sort & Julia name & Grade & Claim \\
\midrule
\endhead
\code{def:typed-sampler}\newline\code{gt-06-types.tex}
& A type graph selects a pair of CL maps sharing one seed & \code{TypedSampler}
& \code{TypedSampler}, \code{sample}, \code{pad_level}
& \constructed\ base & C4b \textsc{Conj.} \\
\code{sec:ld-compiler}\newline\code{gt-10-answer-reduction.tex}
& Six low-degree copies form an 18-type family & \code{PCPCLMap}
& \code{pcp_sampler}, \code{intrinsic_pcp_levels}
& local \grade{CHECKED}; held & C4b \textsc{Conj.} \\
\code{sec:ld-compiler}\newline\code{gt-10-answer-reduction.tex}
& Product with the oracularized sampler has 54 types &
\code{TypedProductCL} & \code{oracularize_sampler},
\code{typed_sampler_product} & \grade{CHECKED} shape; theorem cited
& C4b \textsc{Conj.} \\
\code{fig:decider-pcp}\newline\code{gt-10-answer-reduction.tex}
& Five guarded checks connect input, proof, low degree, and PCP &
\code{TypedAnswer}\newline\code{ReducedDecider} & \code{typed_answer_reduced_decider}
& \grade{CHECKED} finite paths & no current row \\
\code{lem:detyping-verifiers}\newline\code{gt-06-types.tex}
& Remove types at cost \(+2\) levels and \(16^{54}\) soundness factor &
\code{CitedDetyped}\newline\code{Verifier} & \code{detype} & \grade{CITED} & C4b \textsc{Conj.} \\
\code{thm:ar}\newline\code{gt-10-answer-reduction.tex}
& Conditional complexity, completeness, soundness, entanglement contract &
\code{TypedAnswer}\newline\code{ReducedVerifier} & \code{answer_reduce_pcp},
\code{AnswerReduce} & \grade{Checked}/\newline\grade{Cited} & C5 \textsc{Sketch} \\
\bottomrule
\end{longtable}
\normalsize

\subsection{Midpoint diagnostic (TB0.5)}

\scriptsize
\begin{longtable}{P{0.21\textwidth}P{0.16\textwidth}P{0.19\textwidth}P{0.19\textwidth}P{0.12\textwidth}P{0.08\textwidth}}
\toprule
Analytic anchor & Analytic statement & IR sort & Implementation & Grade & Claim \\
\midrule
\endhead
\code{handoff.md}\newline midpoint diagnostic
& False level-\(n\) claim has value \(1-2^{-n}\) under the stated orbit condition
& \code{Test/Ask/Coin} terms & \code{toys/midpoint}
& proof plus exact tests & C6 \textsc{Proved} \\
\code{toys/midpoint/PROOF.md}
& Sequential \(r\)-copy AND value is \((1-2^{-n})^r\); constant gap costs
\(\Theta(2^n)\) copies & exact dynamic program & \code{sequential_value}
& proof plus exact tests & N1 \textsc{Proved} \\
\bottomrule
\end{longtable}
\normalsize

\subsection{Compression and fixed point (planned TB4)}

\scriptsize
\begin{longtable}{P{0.21\textwidth}P{0.16\textwidth}P{0.19\textwidth}P{0.19\textwidth}P{0.12\textwidth}P{0.08\textwidth}}
\toprule
Paper object & Analytic statement & IR sort & Julia name & Grade & Claim \\
\midrule
\endhead
\code{fig:compress}\newline\code{gt-12-compression.tex}
& \(\Repeat\circ\AnswerReduce\circ\Introspect\) &
\code{ClosedProgram}\newline\code{{Compressor}} & \code{Compress} & \grade{CITED}; absent
& C7 \textsc{Conj.} \\
\code{thm:compression}\newline\code{gt-12-compression.tex}
& Nine-level gap-preserving compression with dimension bound & contract tree
& no current implementation & \grade{CITED} & C7 \textsc{Conj.} \\
\code{fig:halt_f}\newline\code{gt-12-compression.tex}
& Self-application, Kleene, and YCode coincide by
Theorems~\ref{thm:kleene-two}--\ref{thm:fixed-point-equivalence} &
\code{PartialProgram}, \code{Quoted} & \code{Hole}, \code{Specialize},
\code{Fix}/\code{YCode} & \grade{CITED}; absent & none \\
\bottomrule
\end{longtable}
\normalsize

\section{What is analytic and what is computed}

\subsection{The evidence boundary}

The analytic content is the sequence of transformations and the implications
attached to them.  The computer currently realizes a strict middle segment.

\begin{center}
\small
\begin{tabularx}{\textwidth}{P{0.20\textwidth}P{0.35\textwidth}X}
\toprule
Class & Present content & What would be required to move the boundary \\
\midrule
Executed constructions
& Circuit/Tseitin syntax; formula arithmetization; multilinear extensions;
\(c_0\); zero-basis quotients; point/line maps; typed products; guarded TB2
decider
& Scale-independent front end for quoted programs and a constructor-faithful
lazy CL stage representation for the PCP sampler \\
Certified identities
& Normalized coefficient equality; zero remainder; structural and actual
degrees; dependency blocks; CL marginals; TB1 exact histograms; finite parser
and branch checks
& General semantic certificates connecting trace, SAT, Tseitin, and
arithmetization, plus promotion through adversarial verdicts \\
Cited analytic leaves
& Cook--Levin asymptotics; general Tseitin semantics; PCP soundness;
low-degree extraction; oracularization; detyping; introspection; rigidity;
anchored repetition; compression
& Formal proofs over quantified machines, distributions, Hilbert spaces,
measurements, approximation metrics, and asymptotic bounds---not more finite
samples \\
\bottomrule
\end{tabularx}
\end{center}

Lean or another proof assistant could check the algebraic core once finite
fields, sparse polynomials, interpolation, and polynomial division are
formalized.  That would move the general zero-basis identity and degree lemmas
from prose to proof terms.  It would not by itself certify
\code{lem:ld-soundness}: that leaf requires a library for finite-dimensional
Hilbert spaces, POVMs, state-dependent distances, tensor-code tests, and the
quantitative extraction argument.  Rigidity and oracularization require the
same operator layer.  Detyping additionally needs a formal probability and
rejection-sampling argument with the \(16^{|\mathrm{Type}|}\) loss.

Anchoring and parallel repetition are combined in \(\Repeat\) and are
theorem-backed rather than executed here.  Introspection and compression would
require a still broader object: a
formal semantics for quoted verifier descriptions, resource bounds under
specialization and universal evaluation, and a proof that the directed value
and entanglement implications compose.  The relevant final statements are
\code{fig:compress} and \code{thm:compression} in
\code{ground-truth/gt-12-compression.tex}.

\subsection{The midpoint diagnostic}

The recursive midpoint verifier for \(y=f^{2^n}(x)\) asks for \(z\) and checks
one of
\[
 z=f^{2^{n-1}}(x),\qquad y=f^{2^{n-1}}(z)
\]
with a fresh fair coin.  It has a compact fixed-point description and perfect
completeness.  Under the orbit-prefix condition stated in C6, the exact optimum
for a false claim is nevertheless
\[
 p_n=1-2^{-n}.
\]
Claim C6 is \textsc{PROVED}.  For adaptive sequential AND repetition, N1 is
also \textsc{PROVED} and gives value \((1-2^{-n})^r\); reaching error at most
\(1/2\) requires \(r=\Theta(2^n)\), hence \(\Theta(n2^n)\) transcript cost in
the stated unit-cost model.  Neither claim concerns quantum parallel repetition
or the actual \(\Compress\) operator.

The diagnostic isolates the point of the whole construction.  A recursive
term can represent an exponentially long computation, and a fixed point can
make that term self-referential, while the acceptance gap collapses too quickly
for iteration.  Representability is easy; an iteratable, gap-preserving
analytic transformation is the hard requirement.

\subsection{Final accounting}

The lambda reformulation removes none of the mathematically difficult leaves.
It makes quotation, hardwiring, description size, timeout, and self-reference
explicit in one syntax.  The polynomial IR makes the Cook--Levin-to-PCP middle
visible as actual algebra, and the CL datatype exposes several sampler closure
laws.  What remains unchanged is the proof that high success in the typed
quantum game yields consistent low-degree polynomial measurements, survives
oracularization and detyping with the stated quantitative loss, survives
introspection and anchored parallel repetition, and composes into the
gap-preserving compression theorem.  Those are analytic theorems, presently
\grade{CITED}; they are neither consequences of homoiconicity nor conclusions
of the finite computations reported here.

